\begin{table}[ht!] 
\centering
\setlength{\tabcolsep}{3pt}
\renewcommand{\arraystretch}{1.4}
\captionsetup{font=small}
\caption{\textbf{Feature comparison of contemporary voice agent evaluation frameworks.} \pmark~denotes partial support, and {---} for the simulator validation means it doesn't apply because of missing simulator. \framework~is the only framework combining live multi-turn simulation across both speech-to-speech (S2S) and cascade architectures with a realistic audio environment, automated simulator validation, comprehensive metrics exposing a wide range of voice agent failures, and pass\^{}k measurement via multi-trial consistency measurement.}
\resizebox{\textwidth}{!}{%
\begin{tabular}{@{}p{1.4in}p{1.5in}p{1.0in}cccccc@{}}
\toprule
\textbf{Framework}
    & \textbf{Interaction}
    & \shortstack[l]{\textbf{Supported}\\\textbf{Architectures}}
    & \shortstack{\textbf{Multi-}\\\textbf{turn}}
    & \shortstack{\textbf{Tool}\\\textbf{Use}}
    & \shortstack{\textbf{Realistic}\\\textbf{Audio}}
    & \shortstack{\textbf{Simulator}\\\textbf{Validation}}
    & \shortstack{\textbf{Comprehensive}\\\textbf{Metrics}} 
    & \shortstack{\textbf{Multi}\\\textbf{Trial}} \\
\midrule
\rowcolor{green!12}
\textbf{\framework}
    & Live bot-to-bot
    & S2S, Cascade
    & \cmark & \cmark & \cmark & \cmark & \cmark & \cmark  \\[2pt]
%\midrule
$\tau$\textbf{-Voice} \cite{ray2026tauvoicebenchmarkingfullduplexvoice}
    & Live bot-to-bot
    & S2S
    & \cmark & \cmark & \cmark & \xmark & \xmark & \xmark \\[2pt]
\textbf{FDB-v3} \cite{lin2026full}
    & Real human audio
    & S2S, Cascade
    & \xmark & \cmark & \pmark & {---} & \xmark & \xmark \\[2pt]
\textbf{VoiceAgentBench}~\cite{jain2025voiceagentbench}
    & Static, TTS-synthesized
    & S2S, Cascade
    & \cmark & \cmark & \xmark & {---} & \xmark & \xmark \\[2pt]
\textbf{CAVA} \cite{cava2025}
    & Partial simulation
    & S2S, Cascade
    & \cmark & \cmark & \xmark & \xmark & \xmark & \xmark \\[2pt]
\textbf{FDB-v2} \cite{lin2026fdb_v2}
    & Live, auto examiner
    & S2S
    & \cmark & \xmark & \xmark & \xmark & \xmark & \xmark \\[2pt]
\textbf{FD-Bench} \cite{peng2025fd}
    & Live, simulated
    & S2S
    & \xmark & \xmark & \xmark & \xmark & \xmark & \xmark \\
\bottomrule
\end{tabular}%
}
\label{tab:framework_comparison}
\end{table}
\section{Related Work}


Many existing voice benchmarks focus on individual components such as STT robustness \citep{chen2024voicebench, ao2024sd, cui2025voxeval}, TTS quality \citep{manku2026emergentttseval, huang2025sheet}, or conversational dynamics \citep{peng2025fd, aroratalking}, rather than the end-to-end behavior of a voice agent. We organize the following discussion around the two challenges introduced above: the fidelity of multi-turn simulation and the comprehensiveness of voice agent quality measurement.


\textbf{Conversation Simulation.} Effective voice agent evaluation requires a simulation methodology that faithfully replicates the dynamic, real-time nature of spoken interaction, where the agent must navigate complete, task-oriented multi-turn conversations with live users whose requests and clarifications may shift throughout the call. Several benchmarks fall short on this requirement in distinct ways. FullDuplex-Bench-v1 (FDB) and FDB-v1.5 \citep{lin2025full, lin2025fullv15} assess conversation dynamics in a heavily-scripted manner without task completion or tool use, rendering themselves unsuitable for voice agent evaluations. VoiceAgentBench \citep{jain2025voiceagentbench} evaluates multi-tool workflows but relies on static TTS-synthesized queries with no conversational back-and-forth. FDB-v3 \citep{lin2026full} improves realism via authentic human recordings with disfluency annotation, yet remains single-turn. Both are further constrained by fixed interactions that limit generalization to unseen scenarios. $\tau$-Voice \citep{ray2026tauvoicebenchmarkingfullduplexvoice} and FDB-v2 \citep{lin2026fdb_v2} represent the closest prior work in terms of live bot-to-bot simulation over multi-turn interactions. However, neither provides automated validation of simulator behavior across trials, leaving open the question of whether evaluation scores reflect agent quality or simulator variance. Furthermore, in $\tau$-Voice, accent variation is coupled with changes in user persona and behavioral style, making it difficult to isolate the acoustic effect of accent from confounding behavioral differences. We address these gaps by introducing a live, multi-trial, bot-to-bot conversation simulator with a controlled perturbation suite and automatic user simulator quality validation.


\textbf{Voice Agent Quality Measurement.} Existing benchmarks that evaluate voice agent behavior converge on a narrow set of metrics. VoiceAgentBench \citep{jain2025voiceagentbench} reports tool selection accuracy and structural consistency of tool invocations, but does not assess any dimension of conversational quality. $\tau$-Voice \citep{ray2026tauvoicebenchmarkingfullduplexvoice} improved on this with a suite of turn-taking measures (response rate, latency, interruption rate, and selectivity) but does not assess whether the agent communicated faithfully or appropriately throughout the interaction. FDB-v3 \citep{lin2026full} introduces a response quality dimension judged at the transcript level and latency decomposition, but does not assess policy faithfulness or accuracy of spoken entities at audio level. To the best of our knowledge, none of these frameworks measure whether the agent makes efficient progress, avoids imposing excessive cognitive load on the user, or speaks the correct information. Collectively, a substantial portion of voice agent quality remains unmeasured, particularly those most consequential for enterprise deployment.